\section{Forme d'intersection}

\subsection{Première approche : Une définition géométrique}
Avant de formaliser la forme d'intersection par le biais du cup-produit sur les classes de cohomologie, il est remarquable d'en donner l'approche géométrique historique. Cette formulation, restreinte aux variétés lisses, repose sur l'intersection transversale de sous-variétés et fournit l'intuition visuelle essentielle pour comprendre les invariants de dimension $4$.

Dans toute cette section, $M$ désigne une variété différentielle fermée, compacte, orientée de dimension $4$.

\begin{definition}[Classe homologique d'un plongement]
Pour une surface fermée orientée lisse $\Sigma$ plongée dans $X$, on notera \emph{par abus de notation} $[\Sigma] \in H_2(X;\ZZ)$ la classe d'homologie correspondante, image de la classe fondamentale $[\Sigma] \in H_2(\Sigma;\ZZ)$ par le morphisme induit par le plongement.
\end{definition}

\begin{theorem}[Représentants lisses]
\label{thm:smooth_rep}
Soit $X$ une 4-variété fermée lisse orientée. Pour toute classe $\alpha \in H_2(X;\ZZ)$, il existe une surface fermée, orientée, lisse et \emph{plongée} $\Sigma \hookrightarrow X$ telle que $[\Sigma] = \alpha$.
\end{theorem}

\begin{proof}

A voir.

\end{proof}

\begin{remark}
En pratique, \textbf{en dimension 4, on peut toujours représenter une classe de $H_2$ par une surface lisse plongée (ou au moins immergée avec des intersections transverses que l'on peut contrôler algébriquement)}.
\end{remark}

\begin{lemma}[Invariance de classe d'homologie sous petite perturbation]
Si $\Sigma_2'$ est une perturbation de $\Sigma_2$ (isotopie), alors $[\Sigma_2'] = [\Sigma_2] \in H_2(X;\ZZ)$.
\end{lemma}
\begin{proof}
Une isotopie $F : \Sigma_2 \times [0,1] \to X$ fournit une 3-chaîne (une variété à bord) dont le bord est $\Sigma_2' - \Sigma_2$. En effet, considérons $W = F(\Sigma_2 \times [0,1])$. C'est une sous-variété de dimension $3$ à bord dans $X$, car $F$ est un plongement. Son bord est
\[
\partial W = F(\Sigma_2 \times \{1\}) \cup F(\Sigma_2 \times \{0\}) = \Sigma_2' \cup \Sigma_2.
\]
En choisissant l'orientation adéquate sur $W$ (celle induite par l'orientation de $\Sigma_2$ et l'orientation de $[0,1]$), le bord orienté est $\partial W = \Sigma_2' - \Sigma_2$ (le signe $-$ provient du fait que l'orientation de $\Sigma_2 \times \{0\}$ est opposée à celle de $\Sigma_2$ pour la convention du bord sortant). Ainsi, $[\Sigma_2'] - [\Sigma_2]$ est un bord dans $C_2(X)$, donc nul en homologie :
\[
[\Sigma_2'] = [\Sigma_2] \in H_2(X;\ZZ).
\]
\end{proof}



\subsubsection{Transversalité}

\begin{definition}[Transversalité de deux sous-variétés]
Soient $M^m$ et $N^n$ deux sous-variétés lisses d'une variété ambiante $X^k$ (avec $m, n \leq k$). On dit que $M$ et $N$ sont \textbf{transverses} en $p \in M \cap N$ si
\[
  T_p M + T_p N = T_p X,
\]
c'est-à-dire si $T_p M$ et $T_p N$ engendrent $T_p X$. On dit que $M$ et $N$ sont \textbf{transverses} (en notation $M \pitchfork N$) si elles sont transverses en tout point de $M \cap N$.

Lorsque $m + n = k$, la condition de transversalité en $p$ est équivalente à $T_p M \cap T_p N = \{0\}$. Dans le cas général, la dimension attendue de $M \cap N$ est $m + n - k$ (éventuellement négative, auquel cas $M \cap N = \emptyset$ génériquement).
\end{definition}

\begin{remark}
Si $M \pitchfork N$, alors $M \cap N$ est une sous-variété lisse de $X$ de dimension $m + n - k$. En particulier :
\begin{itemize}
  \item si $m + n < k$, la transversalité implique $M \cap N = \emptyset$ ;
  \item si $m + n = k$, l'intersection est un ensemble discret (fini si $M$ et $N$ sont compactes) ;
  \item si $m + n > k$, l'intersection est une sous-variété de dimension $m + n - k > 0$.
\end{itemize}
\end{remark}

\begin{definition}[Signe d'une intersection transverse]
Soient $M^m$ et $N^n$ deux sous-variétés orientées plongées transversalement dans $X^{m+n}$ orientée. En un point $p \in M \cap N$, on définit le \textbf{signe d'intersection} $\varepsilon_p \in \{+1, -1\}$ par
\[
  \varepsilon_p = +1 \iff \text{la base orientée de } T_p M \oplus T_p N \text{ est positivement orientée dans } T_p X.
\]
Autrement dit, si $(e_1, \ldots, e_m)$ est une base orientée de $T_pM$ et $(f_1, \ldots, f_n)$ une base orientée de $T_pN$, alors
\[
  \varepsilon_p = \mathrm{sgn}\det(e_1, \ldots, e_m, f_1, \ldots, f_n),
\]
où le déterminant est calculé dans toute base orientée de $T_pX$.
\end{definition}

\begin{remark}
Le signe $\varepsilon_p$ est bien défini (indépendant des bases choisies) et vérifie $\varepsilon_p(M, N) = (-1)^{mn}\, \varepsilon_p(N, M)$, ce qui reflète l'anticommutativité du produit d'intersection en cohomologie.
\end{remark}



\begin{definition}[Nombre d'intersection]
\label{def:intersection_number}
Soient $\Sigma_1, \Sigma_2$ lisses plongés et transverses. On définit
\[
  \Sigma_1 \cdot \Sigma_2 \;=\; \sum_{p \in \Sigma_\alpha \cap \Sigma_\beta} \varepsilon_p \;\in\;\ZZ.
\]
\end{definition}

\begin{definition}[Bonne définition]
\label{lem:welldefined}
On définit $\alpha \cdot \beta$ = $\Sigma_\alpha \cdot \Sigma_\beta$ où $\Sigma_\alpha$ et $\Sigma_\beta$ sont des représentants lisses plongés transverses de $\alpha$ et $\beta$ respectivement (ce qui est possible par le Théorème~\ref{thm:smooth_rep} et le Lemme~\ref{lem:transv}) et l'on a indépendance du choix des représentants transverses.
\end{definition}
\begin{proof}
A voir.
\end{proof}

\begin{definition}[Forme d'intersection]
La \textbf{forme d'intersection} de $X$ est l'application
\[
  Q_X : H_2(X;\ZZ) \times H_2(X;\ZZ) \to \ZZ, \quad Q_X(\alpha,\beta) = \alpha \cdot \beta.
\]

.
\end{definition}




\subsection{Définition homologique classique}

Bien que l'approche par intersection transversale de sous-variétés lisses soit la plus intuitive, elle restreint la définition aux variétés munies d'une structure différentielle. Pour étendre cet invariant aux variétés topologiques quelconques (cadre indispensable pour appréhender les théorèmes de Freedman), il est nécessaire de formaliser la forme d'intersection directement sur les classes d'homologie singulière $H_2(M)$.

Dans toute cette sous-section, $M$ désigne une variété topologique fermée, compacte et orientée de dimension $4$, et $[M] \in H_4(M)$ désigne sa classe fondamentale d'orientation.

La définition purement homologique repose sur l'utilisation combinée du crochet de dualité et de l'isomorphisme de Poincaré. Rappelons que le cap-produit par la classe fondamentale induit un isomorphisme de $\mathbb{Z}$-modules :
\[
D : H^2(M) \longrightarrow H_2(M), \quad \alpha \mapsto [M] \frown \alpha
\]

Puisque cet opérateur est bijectif, on peut inverser cette correspondance. À chaque classe d'homologie de dimension deux $a \in H_2(M)$, on associe de manière unique sa classe de cohomologie duale de Poincaré, notée $D^{-1}(a) \in H^2(M)$.

Nous pouvons désormais poser la définition algébrique classique sur le second groupe d'homologie, libre de toute considération différentielle.

\begin{definition}[Forme d'intersection homologique]
Soit $M$ une variété topologique fermée, orientée de dimension $4$. On appelle \emph{forme d'intersection} de $M$ l'application bilinéaire symétrique entière notée $Q_M$ et définie par :
\[
Q_M : H_2(M)/\text{Tor} \times H_2(M)/\text{Tor} \longrightarrow \mathbb{Z}
\]
\[
Q_M(a, b) = \langle D^{-1}(a) \smile D^{-1}(b), [M] \rangle
\]
où $\smile$ est le cup-produit sur $H^2(M)$ et $\langle \cdot, \cdot \rangle$ désigne le crochet d'évaluation sur la classe fondamentale.
\end{definition}

Le quotient par le sous-module de torsion $\text{Tor}$ est rigoureusement requis ici : si un élément $a$ est de torsion (c'est-à-dire qu'il existe $k \in \mathbb{Z}^*$ tel que $k a = 0$), son produit avec n'importe quelle classe donnerait $0$ par linéarité, ce qui rendrait la forme dégénérée. Le quotient $H_2(M)/\text{Tor}$ est un $\mathbb{Z}$-module libre de rang fini, ce qui permet d'adopter des notations matricielles.


Cette formulation permet de démontrer de manière purement algébrique les propriétés fondamentales du réseau associé à la variété :

\begin{proposition}[Propriétés de structure]
La forme d'intersection $Q_M$ définie sur $H_2(M)/\text{Tor} \simeq \mathbb{Z}^r$ vérifie :
\begin{enumerate}
    \item \textbf{Symétrie :} Pour tout $a, b \in H_2(M)/\text{Tor}$, $Q_M(a, b) = Q_M(b, a)$. Cela découle immédiatement de la graduée-commutativité du cup-produit : pour deux classes de degré $2$, on a $D^{-1}(a) \smile D^{-1}(b) = (-1)^{2 \times 2} D^{-1}(b) \smile D^{-1}(a) = D^{-1}(b) \smile D^{-1}(a)$.
    \item \textbf{Unimodularité :} L'application d'adjointement $\phi : H_2(M)/\text{Tor} \to (H_2(M)/\text{Tor})^*$ induite par $Q_M$ est un isomorphisme de $\mathbb{Z}$-modules. Cela garantit que la matrice représentant $Q_M$ dans n'importe quelle base possède un déterminant égal à $\pm 1$.
\end{enumerate}
\end{proposition}

\begin{remark}[Lien avec le cap-produit]
En utilisant la relation d'adjonction fondamentale liant le cup-produit et le cap-produit, l'expression de la forme d'intersection peut se réécrire sous une forme plus compacte ne faisant intervenir qu'une seule dualité :
\[
Q_M(a, b) = \langle D^{-1}(a), [M] \frown D^{-1}(b) \rangle = \langle D^{-1}(a), b \rangle
\]
Cette écriture montre que calculer la forme d'intersection de deux cycles $a$ et $b$ revient exactement à évaluer la classe de cochaînes duale de $a$ sur le cycle homologique $b$.
\end{remark}



Par la suite, on proposera néanmoins de poursuivre avec la définition géométrique étant plus intuitive. Cette section comporte quelques propriétés intéressantes concernant la forme d'intersection qui nous permettront ensuite de faire le lien avec les résultats de classifications énoncés plus tôt dans le mémoire.

\subsection{Propriétés}

\begin{proposition}[$Q_X$ est symétrique]
\label{prop:sym}
Pour toutes $\alpha, \beta \in H_2(X;\ZZ)$ :
\[
  Q_X(\alpha, \beta) = Q_X(\beta, \alpha).
\]
\end{proposition}

\begin{proof}
Soient $\Sigma_\alpha \pitchfork \Sigma_\beta$. En un point $p \in \Sigma_\alpha \cap \Sigma_\beta$, le signe $\varepsilon_p(\Sigma_\alpha, \Sigma_\beta)$ est $+1$ ssi la base $(e_1, e_2, f_1, f_2)$ est positive dans $T_pX$, où $(e_1,e_2)$ est une base orientée de $T_p\Sigma_\alpha$ et $(f_1,f_2)$ de $T_p\Sigma_\beta$.

Le signe $\varepsilon_p(\Sigma_\beta, \Sigma_\alpha)$ est $+1$ ssi $(f_1, f_2, e_1, e_2)$ est positive dans $T_pX$.

Or la permutation $(e_1,e_2,f_1,f_2) \mapsto (f_1,f_2,e_1,e_2)$ est une permutation de $4$ éléments composée de $4$ transpositions donc signe $(-1)^4 = +1$.

Ainsi $\varepsilon_p(\Sigma_\alpha, \Sigma_\beta) = \varepsilon_p(\Sigma_\beta, \Sigma_\alpha)$, et en sommant sur $p \in \Sigma_\alpha \cap \Sigma_\beta = \Sigma_\beta \cap \Sigma_\alpha$ :
\[
  Q_X(\alpha,\beta) = \sum_p \varepsilon_p(\Sigma_\alpha,\Sigma_\beta) = \sum_p \varepsilon_p(\Sigma_\beta,\Sigma_\alpha) = Q_X(\beta,\alpha). \qedhere
\]
\end{proof}

\begin{remark}
Pour des variétés de dimension congrues à $2 \pmod 4$, la forme d'intersection est \emph{antisymétrique}. La symétrie ici est propre à la dimension 4.
\end{remark}

\begin{proposition}
\label{prop:bilinear}
La forme $Q_X$ est bilinéaire.
\end{proposition}

\begin{proof}
\medskip
\textbf{Linéarité en la première variable.}

Soient $\Sigma, \Sigma'$ des représentants lisses de $\alpha$ et $\alpha'$,
et $T$ un représentant lisse de $\beta$.
 
\medskip
On ne peut pas supposer a priori que $\Sigma$ et $\Sigma'$ sont disjoints dans
$M$.  Cependant, comme une isotopie ambiante induit l'identité en homologie, on
peut remplacer $\Sigma'$ par un représentant isotope $\widetilde\Sigma'$ qui lui
est homologue.
 
\begin{lemma}
Il existe une isotopie ambiante $\phi_t : M \to M$, $t \in [0,1]$, avec
$\phi_0 = \mathrm{id}$, telle que $\widetilde\Sigma' := \phi_1(\Sigma')$
vérifie $\widetilde\Sigma' \cap \Sigma = \emptyset$.
\end{lemma} 
 
On travaille désormais avec $\widetilde\Sigma'$ à la place de $\Sigma'$,
et on suppose de plus $\Sigma$, $\widetilde\Sigma'$, $T$ deux à deux transverses
(position générale).
 
\medskip
Au niveau des chaînes singulières, l'addition en homologie correspond à la
réunion de cycles.  La réunion disjointe $\Sigma \sqcup \widetilde\Sigma'$ est
donc un représentant lisse de $\alpha + \alpha'$ :
\[
  [\Sigma \sqcup \widetilde\Sigma'] \;=\; [\Sigma] + [\widetilde\Sigma']
  \;=\; \alpha + \alpha'.
\]
 
\medskip
Par définition de la forme d'intersection,
\[
  Q(\alpha+\alpha',\,\beta)
  \;=\; \#\bigl((\Sigma \sqcup \widetilde\Sigma') \pitchfork T\bigr).
\]
Comme $\Sigma$ et $\widetilde\Sigma'$ sont disjoints, leurs intersections avec
$T$ le sont aussi :
\[
  (\Sigma \sqcup \widetilde\Sigma') \pitchfork T
  \;=\; (\Sigma \pitchfork T) \sqcup (\widetilde\Sigma' \pitchfork T).
\]
Le signe local $\varepsilon_p = \pm 1$ en chaque point $p$ ne dépend que des
orientations de la surface et de $T$ en $p$, qui n'ont pas été modifiées.
L'additivité du nombre algébrique sur les unions disjointes donne alors
\[
  Q(\alpha+\alpha',\,\beta)
  \;=\; \#(\Sigma \pitchfork T) + \#(\widetilde\Sigma' \pitchfork T)
  \;=\; Q(\alpha,\beta) + Q(\alpha',\beta). \qedhere
\]


\textbf{Homogénéité.}

Soit $\lambda \in \mathbb{Z}$ et $\alpha, \beta \in H_2(M;\mathbb{Z})$.
Un représentant de $\lambda \alpha$ est $|\lambda|$ copies de $\Sigma_\alpha$,
avec orientation renversée si $\lambda < 0$. Par additivité du nombre algébrique
d'intersections sur les réunions disjointes,
\[
  Q(\lambda\alpha, \beta)
  \;=\; \#\bigl((\lambda \cdot \Sigma_\alpha) \pitchfork \Sigma_\beta\bigr)
  \;=\; \lambda \cdot \#(\Sigma_\alpha \pitchfork \Sigma_\beta)
  \;=\; \lambda\, Q(\alpha, \beta). \qedhere
\]

La linéarité en la seconde variable découle, elle, de la symétrie, voir Proposition~\ref{prop:sym}.
\end{proof}

\begin{remark}
L'unimodularité est équivalente au fait que le morphisme adjoint $\varphi_Q : H_2(X;\ZZ) \to \mathrm{Hom}(H_2(X;\ZZ), \ZZ)$ défini par $\varphi_Q(\alpha)(\beta) = Q(\alpha,\beta)$ est un isomorphisme.
\end{remark}

\begin{theorem}[Unimodularité de $Q_M$]
\label{thm:unimod}
Soit $M$ une 4-variété fermée compacte orientée lisse. La forme d'intersection $Q_M$ induit une forme bilinéaire symétrique unimodulaire sur le quotient libre $H_2(M;\mathbb{Z}) / \mathrm{Tor}(H_2(M;\mathbb{Z}))$.
\end{theorem}

\begin{proof}
On observe d'abord que $Q_M$ s'annule sur $\mathrm{Tor}(H_2(M;\mathbb{Z}))$ : si $\alpha \in H_2(M;\mathbb{Z})$ vérifie $n\alpha = 0$ pour un certain $n \geq 1$, alors pour tout $\beta$,
\[
  n \cdot Q_M(\alpha, \beta) = Q_M(n\alpha, \beta) = 0,
\]
et comme $Q_M(\alpha,\beta) \in \mathbb{Z}$, on conclut $Q_M(\alpha,\beta) = 0$. Ainsi $Q_M$ descend bien en une forme bilinéaire sur $H_2(M;\mathbb{Z})_{\mathrm{free}}$, qui est un $\mathbb{Z}$-module libre de type fini.

\medskip

\noindent On veut montrer que
\[
  \Phi : H_2(M;\mathbb{Z})_{\mathrm{free}} \longrightarrow \mathrm{Hom}(H_2(M;\mathbb{Z})_{\mathrm{free}},\mathbb{Z}),
  \qquad \alpha \longmapsto Q_M(\alpha,-)
\]
est un isomorphisme. On écrit $\Phi$ comme la composée
\[
  \Phi \;:\; H_2(M;\mathbb{Z})_{\mathrm{free}}
  \xrightarrow{\;\mathrm{PD}\;}
  H^2(M;\mathbb{Z})
  \xrightarrow{\;\mathrm{ev}\;}
  \mathrm{Hom}(H_2(M;\mathbb{Z})_{\mathrm{free}},\mathbb{Z}),
\]
où $\mathrm{PD}$ est l'isomorphisme de Poincaré et $\mathrm{ev}$ est l'évaluation $\varphi \mapsto (\beta \mapsto \varphi(\beta))$.

\medskip

\noindent Le théorème des coefficients universels donne la suite exacte
\[
  0 \;\longrightarrow\;
  \mathrm{Ext}^1(H_1(M;\mathbb{Z}),\mathbb{Z})
  \;\longrightarrow\;
  H^2(M;\mathbb{Z})
  \xrightarrow{\;\mathrm{ev}\;}
  \mathrm{Hom}(H_2(M;\mathbb{Z}),\mathbb{Z})
  \;\longrightarrow\; 0.
\]
Or $\mathrm{Hom}(H_2(M;\mathbb{Z}), \mathbb{Z}) = \mathrm{Hom}(H_2(M;\mathbb{Z})_{\mathrm{free}}, \mathbb{Z})$ car tout homomorphisme vers $\mathbb{Z}$ s'annule sur la torsion. De plus, par la dualité de Poincaré, $H_2(M;\mathbb{Z})_{\mathrm{free}} \cong H^2(M;\mathbb{Z}) / \mathrm{Ext}^1(H_1(M;\mathbb{Z}), \mathbb{Z})$, de sorte que $\mathrm{ev}$ induit un isomorphisme
\[
H^2(M;\mathbb{Z}) \;\xrightarrow{\;\sim\;}\; \mathrm{Hom}(H_2(M;\mathbb{Z})_{\mathrm{free}}, \mathbb{Z}).
\]

\medskip

\noindent Ainsi $\Phi = \mathrm{ev} \circ \mathrm{PD}$ est une composée de deux isomorphismes, donc un isomorphisme, ce qui conclut.
\end{proof}

%==============================================================================
\subsection{Exemples}
%==============================================================================

\begin{exemple}[$S^4$]
On a $H_2(S^4;\mathbb{Z}) = 0$, donc $Q$ est la forme nulle sur le groupe
trivial.
\end{exemple}

\begin{exemple}[$S^2 \times S^2$]
On a $H_2(S^2 \times S^2;\mathbb{Z}) = \mathbb{Z}^2$, engendré par
\[
  \alpha = [S^2 \times \{p\}], \qquad \beta = [\{p\} \times S^2]
\]
pour un point $p \in S^2$ quelconque. On munit $S^2 \times S^2$ de
l'orientation produit, et chaque facteur de son orientation standard.

\medskip
\noindent\textbf{Calcul de $Q(\alpha,\alpha)$.}
Pour calculer $Q(\alpha,\alpha)$, on prend deux représentants distincts :
$\Sigma = S^2 \times \{p\}$ et $\Sigma' = S^2 \times \{q\}$ avec $p \neq q$.
Ces deux surfaces sont disjointes dans $S^2 \times S^2$, donc
\[
  Q(\alpha, \alpha) = \#(\Sigma \pitchfork \Sigma') = 0.
\]
De même $Q(\beta,\beta) = 0$.

\medskip
\noindent\textbf{Calcul de $Q(\alpha,\beta)$.}
On prend $\Sigma = S^2 \times \{p\}$ et $T = \{p\} \times S^2$. Ces deux
surfaces se coupent transversalement en l'unique point $(p,p)$. Il reste à
déterminer le signe $\varepsilon(p,p) \in \{+1,-1\}$.

Par définition, $\varepsilon(p,p) = +1$ si et seulement si la concaténation
des bases orientées
\[
  (T_{(p,p)}\Sigma,\; T_{(p,p)}T)
  \;=\;
  (T_p S^2 \times \{0\},\; \{0\} \times T_p S^2)
\]
est une base positivement orientée de $T_{(p,p)}(S^2 \times S^2) \cong T_p S^2
\oplus T_p S^2$. Or l'orientation produit de $S^2 \times S^2$ est précisément
celle définie par cette décomposition, donc la concaténation est bien une base
positive. Ainsi $\varepsilon(p,p) = +1$ et
\[
  Q(\alpha,\beta) = 1.
\]
La matrice de $Q$ dans la base $(\alpha,\beta)$ est donc
\[
  Q_{S^2 \times S^2} \;=\; \begin{pmatrix} 0 & 1 \\ 1 & 0 \end{pmatrix}.
\]
\end{exemple}

\begin{exemple}[$\mathbb{CP}^2$]
On munit $\mathbb{CP}^2$ de son orientation complexe canonique (celle pour
laquelle les cartes holomorphes sont directement orientées). On a
$H_2(\mathbb{CP}^2;\mathbb{Z}) = \mathbb{Z}$, engendré par $h = [\mathbb{CP}^1]$
la classe d'une droite projective complexe, munie de son orientation complexe.

\medskip
\noindent\textbf{Calcul de $Q(h,h)$.}
On prend deux droites projectives distinctes en position générale :
\[
  L_1 = \{[z_0 : z_1 : 0]\} \subset \mathbb{CP}^2, \qquad
  L_2 = \{[z_0 : 0 : z_2]\} \subset \mathbb{CP}^2.
\]
Ces deux droites se coupent transversalement en l'unique point $[1:0:0]$. Il
reste à calculer le signe en ce point.

En coordonnées affines $z_0 = 1$, on paramètre $L_1$ par $(z_1, 0)$ et $L_2$
par $(0, z_2)$, avec $z_1, z_2 \in \mathbb{C}$. Les espaces tangents réels en
$[1:0:0]$ sont
\[
  T_{[1:0:0]} L_1 = \mathbb{R}\,\partial_{x_1} \oplus \mathbb{R}\,\partial_{y_1},
  \qquad
  T_{[1:0:0]} L_2 = \mathbb{R}\,\partial_{x_2} \oplus \mathbb{R}\,\partial_{y_2}
\]
où $z_k = x_k + iy_k$. L'orientation complexe de $L_k \cong \mathbb{C}$ est
$(\partial_{x_k}, \partial_{y_k})$. La concaténation
\[
  (\partial_{x_1}, \partial_{y_1}, \partial_{x_2}, \partial_{y_2})
\]
est une base de $T_{[1:0:0]}\mathbb{CP}^2 \cong \mathbb{C}^2$. L'orientation
complexe canonique de $\mathbb{CP}^2$ est précisément celle définie par la base
$(\partial_{x_1}, \partial_{y_1}, \partial_{x_2}, \partial_{y_2})$, donc le
signe est $+1$. Ainsi
\[
  Q(h,h) = 1.
\]
La matrice de $Q$ dans la base $(h)$ est donc
\[
  Q_{\mathbb{CP}^2} \;=\; \begin{pmatrix} 1 \end{pmatrix}.
\]
\end{exemple}

%==============================================================================
\subsection{Invariants topologiques via la forme d'intersection}
%==============================================================================

Les invariants algébriques classiques des formes bilinéaires symétriques entières, introduits dans la section [?], trouvent ici une interprétation topologique directe lorsqu'on les applique à la forme d'intersection $Q_X$ d'une $4$-variété $X$. 

\subsubsection{Rang et nombres de Betti}

Puisque le $\mathbb{Z}$-module libre $H_2(X;\mathbb{Z}) / \operatorname{Tors}$ est de rang fini, le rang de la forme d'intersection $Q_X$ est bien défini. La dualité de Poincaré garantit que $Q_X$ est unimodulaire, et donc en particulier non dégénérée. Par conséquent, son rang coïncide exactement avec le deuxième nombre de Betti de la variété :
\[
\operatorname{rg}(Q_X) = b_2(X).
\]

\subsubsection{Signature}

De même, la signature $\sigma(Q_X) = b_2^+ - b_2^-$ définie précédemment devient un invariant fondamental de la $4$-variété, que l'on note simplement $\sigma(X)$. Les dimensions $b_2^+$ et $b_2^-$ des sous-espaces maximaux sur lesquels $Q_X \otimes \mathbb{R}$ est définie positive (resp. négative) caractérisent le type de la variété : on dit que $X$ est \emph{définie} (positive ou négative) ou \emph{indéfinie} selon la nature algebrique de sa forme.

La signature est particulièrement robuste vis-à-vis des opérations chirurgicales et topologiques sur les variétés. Notamment, le renversement d'orientation inversant le signe de la forme, on a immédiatement :
\[
\sigma(-X) = -\sigma(X).
\]
De plus, la structure additive de l'homologie par rapport à la somme connexe implique que la forme d'intersection d'une somme connexe est la somme directe des formes de chaque facteur. La signature étant un invariant additif pour la somme directe, on obtient :
\[
\sigma(X_1 \mathbin{\#} X_2) = \sigma(X_1) + \sigma(X_2).
\]
Enfin, la signature est une obstruction au cobordisme orienté :

\begin{lemma}
Si $M^4$ est le bord d'une $5$-variété orientée compacte $W^5$, alors $\sigma(Q_M) = 0$.
\end{lemma}

La réciproque de ce lemme est fausse en général : la nullité de la signature ne suffit pas à garantir que $M^4$ borde une $5$-variété, d'autres obstructions pouvant intervenir. En revanche, dans le cadre lisse, Rokhlin a démontré la réciproque suivante, bien plus difficile.

\begin{theorem}[Rokhlin]
Si $M^4$ est une $4$-variété lisse orientée avec $\sigma(Q_M) = 0$, alors il existe une $5$-variété lisse orientée $W^5$ telle que $\partial W = M$.
\end{theorem}
\newpage

%========================================================================

